\chapter{Analyse van de context}
\label{ch:contextanalyse}

Dit hoofdstuk beschrijft de organisatorische en technische context waarbinnen impactanalyses bij VLAIO plaatsvinden op het moment van de start van dit onderzoek. Eerst wordt de data-architectuur geschetst, gevolgd door de huidige mogelijkheden en beperkingen van lineage binnen Databricks en Unity Catalog. Vervolgens wordt toegelicht hoe impactanalyses vandaag verlopen en welke rol zicht op dataherkomst en afhankelijkheden daarin speelt. Tot slot wordt beschreven hoe OpenMetadata als mogelijke oplossing wordt ingezet, met specifieke aandacht voor de ontbrekende schakel tussen Databricks en Power BI. Dit hoofdstuk vormt daarmee de brug tussen het theoretisch kader uit de literatuurstudie en de proof-of-concept die in het volgende hoofdstuk wordt uitgewerkt.\footnote{Dit hoofdstuk is gebaseerd op interne documentatie en afstemmomenten met het CC-Data team.}

\section{Datalakehouse en medallion-architectuur bij VLAIO}

De data-omgeving van VLAIO is opgebouwd als een grootschalig data lakehouse. Databricks werkt als een laag bovenop die omgeving en is een centraal platform voor toepassingen als data ingestie, verwerking, machine learning en rapportering. In tegenstelling tot een architectuur waarbij een data lake en een datawarehouse apart naast elkaar bestaan, combineert een lakehouse de flexibiliteit van een data lake met de betrouwbaarheid en governance-mogelijkheden van een warehouse. Binnen VLAIO betekent dit concreet dat uiteenlopende bronsystemen hun data aanleveren de Databricks omgeving, waar deze vervolgens via verschillende lagen worden verwerkt tot rapporteerbare datasets.\\

De architectuur volgt het medallion-patroon met opeenvolgende bronze, silver en gold of platinum-lagen. Hierbij worden in de bronze-laag ruwe extracten uit operationele toepassingen opgeslagen, zonder iets van bewerkingen. In de silver-laag worden deze ruwe gegevens opgeschoond, getransformeerd en gestandaardiseerd, bijvoorbeeld door datakwaliteitscontroles of het harmoniseren van velden. De pla\-ti\-num-laag tenslotte bevat businessgerichte, samengevoegde datasets die rechtstreeks gebruikt worden voor rapportering en analyses. Hierop worden dan bijvoorbeeld de Power BI rapporten gebouwd. Door de groei van het aantal bronnen, transformaties en rapporten is deze omgeving in de loop der tijd steeds complexer geworden, met talloze onderlinge afhankelijkheden tussen tabellen, modellen en rapportages.\\

In zo'n context-intensieve omgeving wordt het voor individuele medewerkers onhaalbaar om op basis van geheugen of losse documentatie een volledig overzicht te bewaren van de dataflows. In de literatuur rond metadata management en governance werd reeds vermeld dat in dergelijke complexe omgevingen een gestructureerde aanpak rond data lineage cruciaal is om transparantie, herleidbaarheid en foutopsporing te ondersteunen. De complexiteit van het lakehouse bij VLAIO creëert dus een duidelijke nood aan betrouwbare lineage, zowel om de impact van wijzigingen transparanter en beter herleidbaar te maken, als om de interne governance en controleerbaarheid van rapporten te versterken.\\

\section{Basis-lineage in Databricks en Unity Catalog}

Databricks en de bijhorende Unity Catalog\footnote{Unity Catalog is de ingebouwde governance-laag van Databricks die toegangsbeheer, metadata en lineage centraal beheert over meerdere workspaces heen.} voorzien in principe al een vorm van technische lineage, waarbij op basis van querygeschiedenis en metadata kan worden afgeleid welke tabellen door welke bewerkingen met elkaar verbonden zijn. In de praktijk wordt deze functionaliteit binnen VLAIO slechts beperkt ingezet. Tijdens afstemmomenten werd aangegeven dat lineage in Databricks “niet beschikbaar of niet gemakkelijk beschikbaar” is voor de doorsnee gebruiker, en dat er bovendien een tijdslimiet zit op de historiek die wordt bijgehouden. Daardoor is de zichtbaarheid van afhankelijkheden in de Databricks-omgeving onvolledig en sterk afhankelijk van wie weet waar de lineage-view zich bevindt en hoe deze juist moet worden gebruikt.\\

Daarnaast is deze lineage beperkt tot enkel de data, die in de Databricks omgeving zelf zit, waardoor afhankelijk\-heden met externe applicaties zoals bijvoorbeeld Power BI buiten het beeld blijven. In de besprekingen werd expliciet benoemd dat vooral de sprong van Databricks naar Power BI het moeilijkst te reconstrueren is, onder meer door naamswijzigingen, samenvoegingen en complexe transformaties. Zonder een geïntegreerde, end-to-end lineage moeten medewerkers terugvallen op handmatig zoeken in notebooks en code, wat tijdrovend en foutgevoelig is. Dat maakt de ingebouwde Databricks-lineage weliswaar nuttig als bouwsteen, maar ontoereikend als enige bron voor impactanalyses op organisatieniveau.\\

\section{Impactanalyses bij VLAIO}

Impactanalyses spelen een centrale rol in het datagedreven werken bij VLAIO. In essentie wordt in een impactanalyse onderzocht welke gevolgen een wijziging in databronnen, datamodellen of transformatieprocessen heeft op datasets en rapporten. Typische triggers voor het opstellen van zo een analyse zijn bijvoorbeeld aangekondigde wijzigingen in bronsystemen (velden die verdwijnen of een andere betekenis krijgen) en interne initiatieven om data\-modellen te herontwerpen, bijvoorbeeld om rapportering te vereenvoudigen of nieuwe businessvereisten te ondersteunen.\\

Voor elk geval is het doel om vooraf in kaart te brengen welke datasets, berekeningen en rapporten afhankelijk zijn van de betrokken velden of tabellen, en welke aanpassingen gepland moeten worden voordat de wijziging effectief wordt doorgevoerd. De resultaten van zo'n analyse worden nadien vertaald naar concrete acties voor data-engineers, analisten en Power BI-ontwikkelaars. De huidige impact\-analyse documenten leggen deze redenering vast en beschrijven vaak een as-is en to-be situatie, maar steunen voor de onderliggende reconstructie van datastromen vooral op manuele reconstructie en mondelinge domeinkennis.\\

Het nut van impactanalyses hangt daarmee sterk af van het zicht op onderliggende afhankelijkheden. Zonder overzicht van welke databronnen, tabellen en rapporten door een wijziging geraakt kunnen worden, is het moeilijk om impact betrouwbaar en efficiënt in te schatten. In de huidige situatie vormt het ontbreken van een centraal overzicht van deze afhankelijkheden dus een belangrijke beperking.\\

Voor technische impactanalyses ligt de meerwaarde van lineage vooral in het verminderen van manueel opzoekwerk en het verkleinen van de afhankelijkheid van impliciete kennis bij collega's. In deze bachelorproef wordt dit effect niet afzonderlijk gemeten, aangezien de technische documenten en bijhorende analysepraktijk slechts door een beperkte groep worden gebruikt. Bovendien is de meerwaarde van lineage in die context relatief evident. Een expliciet overzicht van afhankelijkheden maakt het eenvoudiger om betrokken datasets en rapporten te identificeren. Daarom richt deze studie zich naar een breder doelpubliek.\\

De evaluatie in Hoofdstuk~\ref{ch:poc-evaluatie} onderzoekt daarom in welke mate het lineage-overzicht ook voor een ruimere groep rapportgebruikers, bouwers en technische rollen bijdraagt aan het beter begrijpen en inschatten van de impact van wijzigingen op rapporten. Daarbij wordt nagegaan in welke mate extra inzicht in dataherkomst en afhankelijkheden leidt tot meer begrijpbaarheid, herleidbaarheid en vertrouwen in gerapporteerde cijfers.

\section{Huidige pijnpunten in het impactanalyseproces}

De grootste pijnpunten in de huidige situatie bij VLAIO hebben te maken met het ontbreken van een centrale, betrouwbare lineage-view over de volledige keten van bronsystemen tot Power BI-rapporten. Zonder lineage moeten analisten, rapportbouwers en data-engineers telkens opnieuw velden en tabellen manueel volgen doorheen de verschillende data lagen. Dat gebeurt onder meer door notebooks en code door te zoeken, bestandsnamen en kolomnamen te matchen, en collega’s te raadplegen die mogelijk weten waar bepaalde velden nog gebruikt worden. Tijdens afstemmomenten werd vermeld dat men een veld regelmatig ter sprake brengt binnen het team om na te gaan of iemand weet of en waar het gebruikt wordt, wat illustreert hoe sterk het huidige zicht op dataherkomst afhangt van individuele kennis en communicatie.\\

De moeilijkste afhankelijkheden om te achterhalen zijn precies degene waar de standaard tools geen volledig beeld geven, zoals de sprong van Databricks naar Power BI. Daardoor is het niet alleen tijdrovend om de volledige keten te reconstrueren in impactanalyses, maar blijft het ook voor rapportgebruikers moeilijk om bij wijzigingen snel te begrijpen welke cijfers en rapporten geraakt worden en op welke bronnen ze precies steunen.\\

Samengevat leidt de huidige situatie tot lange doorlooptijden bij wijzigingen, een reëel risico op menselijke fouten en een sterke afhanke\-lijk\-heid van specifieke personen met bepaalde domeinkennis. Er is vandaag de dag geen centraal overzicht beschikbaar dat inzicht geeft in de verbindingen tussen databronnen, tabellen en rapporten. Dit terwijl die informatie juist cruciaal is voor zowel impactanalyses, als het begrijpen van rapporten. De POC in deze bachelorproef zal dan ook vanuit dit gebrek vertrekken.

\section{Huidige basis-lineage en de kloof met Power BI}

Om deze knelpunten aan te pakken, wordt binnen VLAIO gewerkt met OpenMetadata als centraal metadata platform. In de literatuurstudie is OpenMetadata gepositioneerd als een open-source oplossing die data cataloging, governance en lineage combineert in één centraal platform. Binnen de context van deze bachelorproef wordt OpenMetadata gebruikt om metadata uit de Databricks-omgeving en andere relevante systemen te centraliseren en om een beter zicht te krijgen op de dataflows in het lakehouse. De OpenMetadata omgeving was reeds opgezet en gekoppeld aan de Databricks service. Via deze connectie werden de tabel en dataset-entiteiten vanuit Databricks reeds geïngesteerd, inclusief de bijbehorende lineage. Hierdoor was alvast een groot deel van de huidige fragmentatie verminderd.\\

Een belangrijke beperking is echter dat de lineage tussen platinum-datasets en de bijbehorende Power BI-rapporten nog niet aanwezig is. Hoewel de Databricks-kant in beeld komt via OpenMetadata, blijft de laatste stap naar Power BI daarmee voorlopig nog buiten het lineagebeeld. In de gesprekken over impactanalyses werd net deze sprong expliciet als problematisch benoemd. Gebruikers willen snel kunnen zien welke rapporten een bepaalde dataset gebruiken en of dezelfde gegevens in verschillende rapporten onderliggend op dezelfde bron steunen. Zonder geïntegreerde Power BI-lineage blijft dit een manueel en foutgevoelig zoekproces, zowel bij impactanalyses als bij het beantwoorden van inhoudelijke vragen over rapporten.\\

In het volgende hoofdstuk wordt onderzocht hoe platinum datasets systematisch aan de bijhorende Power BI rapporten gekoppeld kunnen worden in OpenMetadata. Daarbij zal de logica opgezet worden als een generieke en herbruikbare toepassing, dat gebruikt kan worden voor het volledige Power BI-landschap te verbinden met lineage. Hoofdstuk~\ref{ch:poc-evaluatie} evalueert vervolgens in welke mate het gerealiseerde lineage-overzicht gebruikers effectief helpt om wijzigingen en hun impact beter te begrijpen.